\subsubsection*{Changing the frame at the selected point}
\label{subsec:split-selmer}
\hypertarget{source-split-selmer}{}

To identify those abstract block lengths with the prescribed local conditions, choose integral coordinates adapted to the decomposition at the good point. The surrounding family still includes the points off the theta locus that enter its congruence ideal.

Use the integral family \(R_x\) at the fixed Hecke
localization, including the points outside the theta subset.
At the good theta point \(\pi\) the given rational decomposition
is an isomorphism with \(\rho_\pi\oplus\bigoplus_{k=1}^t\chi_k\).
Here \(\rho_\pi\) is the specified integral rank-\(r\) representation
and \(\chi_k/\chi_1=\epsilon^{1-k}\) for the cyclotomic character
\(\epsilon\). Replace \(R_\pi\) by the direct sum in these
integral bases and keep all other matrices at their own points.
\begin{lemma}[Change of frame in the family algebra]
The rational change of basis identifies the two integral matrix
spans, their marked corners and congruence ideals. At the chosen point \(\pi\), evaluation on the complement gives
\[
 \operatorname{ev}_{\pi,W}(\mathcal A)=M_r(O),
\]
where \(\mathcal A\) is the normalized integral span.
\end{lemma}
\begin{proof}
If \(C_x\) changes the frame at \(x\), conjugation sends each generator
\(R_x(g)\) to \(R'_x(g)\). Thus it carries
\[
 \left(\sum_j a_j(x)R_x(g_j)\right)_x
 \longmapsto\left(\sum_j a_j(x)R'_x(g_j)\right)_x .
\]
The coefficients are the very same Hecke functions. Applying
inverse conjugation proves equality of the integral spans.
Characteristic coefficients are unchanged. Row embeddings and
column quotients transport by rational conjugation, with the
specified lattices of \(\rho_\pi\) and its twisted dual as targets.
The resulting finite family is continuous at each point.
\label{cor:split-selmer-1}
The marked data before the change of frame consist of
orthogonal summand projections at \(\pi\), extended as integral
idempotents in the family span, and a corner matrix \(X_0\) with
distinct roots in every fibre. The fibrewise ranks and the two specializations are
\[
 \operatorname{rank}e_B=t,\qquad \operatorname{rank}e_1=1.
\]

At \(\pi\) the first projects onto the character block and the second onto a line of \(\rho_\pi\). These data can be
transported from an integral marked model on the same coefficient
localization. Write \(R_x^{\rm m}\) for the marked model at \(x\).
The shared Hecke character gives
\[
 \det(1-TR_x(\operatorname{Frob}_v))
   =\det(1-TR_x^{\rm m}(\operatorname{Frob}_v))
\]
at the unramified Frobenius classes. Taking traces, then using
continuity and Chebotarev density, gives equality of the trace
functions on the entire Galois group. Brauer--Nesbitt applied
to the two semisimple rational representations gives an
intertwiner at every \(x\). At \(\pi\), cancelling the specified
character summands identifies the rank-\(r\) complements.
Write \(C\) for the resulting tuple of changes of frame and
\(\mathcal A^{\rm m},\mathcal A'\) for the two integral spans.
The generator identity above gives
\begin{equation}\label{eq:selmer-transport-corner}
 C\mathcal A^{\rm m} C^{-1}=\mathcal A',\qquad
 C(e\mathcal A^{\rm m}e)C^{-1}
      =(CeC^{-1})\mathcal A'(CeC^{-1})
      \quad(e=e_B,e_1).
\end{equation}
By \eqref{eq:selmer-transport-corner} the corner matrix is conjugated in every fibre. Its
characteristic roots, the two idempotent ranks and their
orthogonality are preserved.  \label{cor:split-selmer-2}

The cyclotomic test
\(g_0\) lies in the prescribed inertia and has
\[
 \epsilon(g_0)=1+\ell^s u,
\]
with \(u\) a unit. In the new basis
\(X_0(\pi)\) is \(e_B R_\pi(g_0)e_B\). This identification follows
because both models of the direct sum transport the action of
every sum \(\sum_j a_j(\pi)R_\pi(g_j)\) to the same block action.
Thus the perturbation \(X_0-e_B R(g_0)e_B\) lies in the new span,
is supported on the corner, and vanishes at \(\pi\). The ranks,
corner characteristic polynomials and orthogonality relations
are unchanged by conjugation. The matrix-span and ordered-root
construction consequently applies with the same Hecke localization
and the same marked point.

Let
\[
 B_\pi\subseteq M_r(O)
\]
be the image of the normalized algebra
on the specified rank-\(r\) lattice. Absolute irreducibility of its
reduction gives the full matrix algebra over the residue field by
the density theorem. Writing \(\varpi\) for a uniformizer, this says
\[
 B_\pi+\varpi M_r(O)=M_r(O),\qquad
 M_r(O)/B_\pi=\varpi\bigl(M_r(O)/B_\pi\bigr).
\]
Nakayama's lemma applies to this finite \(O\)-module. Thus \(B_\pi=M_r(O)\): the localized order contains every integral endomorphism of the complementary lattice used in the row calculation.
Section~\ref{subsec:retained-galois} constructs the marked model.
Proposition~\ref{prop:retained-rank-complement} proves residual
absolute irreducibility at this localization.

\end{proof}
\subsubsection*{The lattice used in the rank-\(r\) row quotient}

The new coordinates identify the residual complementary block. Computing its row quotient requires the particular lattice on which its normalized Galois action operates, together with the dual action at the conjugate place.

Write \(U_g=\chi_1(g)^{-1}R(g)\), \(Y=\sum_g\mathcal A(U_g-1)E\),
and \(y(g)=(U_g-1)E\) in the normalized algebra. Let \(T\) be
the specified lattice of the \(U\)-action on the rank-\(r\) block,
so \(W=T[1/\ell]\). Write \(T^*=\operatorname{Hom}_O(T,O)(1)\)
and \(W^*=T^*[1/\ell]\). At a row place the dual finite-class map uses the source lattice and right action
\[
 P=e_1\mathcal A,\qquad g\cdot b=\epsilon(g)bU_g^{-1}.
\]
To see its target, write
\(e_1(\pi)=v\lambda\) with \(\lambda(v)=1\) over \(O\) in the
rank-\(r\) block. An element \(b\in P\) has evaluation \(v\lambda_b\)
on this block. The map \(b\mapsto\lambda_b\) onto \(T^*\) is
surjective because the normalized algebra maps onto \(M_r(O)\).
\begin{lemma}[The local row quotient]
Evaluation \(p:P\to T^*\), \(b\mapsto\lambda_b\), is an integral
surjection for the twisted dual action. With
\(A=W/T,\ L_w=H^1_f(E_w,T^*)\) and
\(J_w=p_*H^1_f(E_w,P)\), put
\(F_w=\operatorname{im}(H^1_f(E_w,W)\to H^1(E_w,A))\).
Local Tate duality induces
\[
 J_w^\perp/F_w\simeq\operatorname{Hom}_O(L_w/J_w,F/O).
\]
\end{lemma}
\begin{proof}
Since the block maps onto \(M_r(O)\), every row \(\lambda_b\) occurs.
If \(\lambda_b\) is a row functional, transposition converts its right action into the column-coordinate action
\[
 \lambda_b^{\mathsf T}\longmapsto\epsilon(g)W(g)^{-\mathsf T}\lambda_b^{\mathsf T}.
\]
This is precisely the action on \(W^*\). For the standard basis \(e_i\), integrality is the coordinate test
\[
 u(O^r)\subseteq O\quad\Longleftrightarrow\quad
 (u(e_1),\ldots,u(e_r))\in O^r
 \qquad(u\in\operatorname{Hom}_F(F^r,F)).
\] Restriction through the
chosen completion identifies the decomposition group with the
local absolute Galois group. The coefficient isomorphism there
carries the fraction field and DVR simultaneously, so it transports
continuous cocycles, coboundaries and the finite conditions with
the lattice. Recording matrix rows at the ordered points embeds the rational source in a finite sum of the evaluated twisted duals. This is the row-colength map.  \label{cor:split-selmer-3}

Use local Tate duality in its \(F/O\)-valued form, obtained from
the usual \(\mathbb Q_\ell/\mathbb Z_\ell\)-valued pairing by the
trace and a generator of the inverse different. An integral class
annihilates \(F_w\) precisely when its rational image pairs to zero
with \(H^1_f(E_w,W)\): any nonzero rational value can be scaled to
a nonintegral one. By Bloch--Kato orthogonality this condition
cuts out \(L_w\). Double annihilation gives
\[
 F_w=L_w^\perp.
\]

The restriction
\[
 J_w^\perp\to\operatorname{Hom}_O(L_w/J_w,F/O)
\]
has this kernel and is onto by injectivity of \(F/O\) over the DVR.
The extension is continuous on the finitely generated integral
cohomology group. Thus the normalized length equals that of
\(L_w/J_w\), bounded by the row index at this same place.
\cite{bloch-kato}
\end{proof}

\begin{lemma}[Annihilation of the integral row classes]
For an \(\mathcal A\)-linear \(\varphi:Y\to A\), put \(c(g)=\varphi(y(g))\). Local Tate duality gives
\[
 \operatorname{res}_w[c]\in J_w^\perp.
\]
\end{lemma}
\begin{proof}
To prove annihilation of \(J_w\), let \(a\) be a continuous
row cocycle with values in \(P\), and let \(c(g)=\varphi(y(g))\)
with values in \(A\). The multiplication cochain
\[ z(g,h)=a(g)U_g y(h) \]
takes values in \(e_1Y\), with the action on this lattice
given by the scalar \(\epsilon\).

The contraction \(z\) is a continuous two-cocycle. This follows by the cocycle identity for \(a\)
with action \(b\mapsto\epsilon(g)bU_g^{-1}\) and the
identity \(y(gh)=y(g)+U_g y(h)\). Its image in
\(e_1\mathcal A E(1)\) is a coboundary. Indeed the
one-cochain \(d(g)=a(g)U_g E\) satisfies
\[
 \epsilon(g)d(h)-d(gh)+d(g)=-a(g)U_g(U_h-1)E=-z(g,h).
\]
The inclusion \(e_1Y\subseteq e_1\mathcal A E\) induces an
injection on continuous \(H^2\) after the Tate twist. Both
lattices are finite free over \(O\), and the local invariant
map
\[
 H^2(E_w,\mathbb Z_\ell(1))\simeq\mathbb Z_\ell,
\]
extended to finite free coefficient modules, identifies the
induced map with the original lattice inclusion. This works
also for a nonsaturated inclusion. Therefore the cochain
\(z\) is already a coboundary in \(e_1Y(1)\).

Apply to that smaller lattice the map \(m\mapsto\lambda(\varphi(m))\).
The pairing is well defined modulo \(O\), and is equivariant
for the scalar Tate action. Since \(\varphi\) is
\(\mathcal A\)-linear, the image of \(z(g,h)\) is
\(\langle p(a(g)),U_g c(h)\rangle\), with
\[
 p:P\to T^*
\]
the row quotient.

This is the cup pairing in the reversed
order, which differs from the desired order by the sign
for two degree-one classes. It follows that \([c]|_{E_w}\)
annihilates the image of every such \([a]\), in particular
\(J_w\). The local class of \(c\) is continuous, as the
universal cocycle is continuous in its finite integral
lattice and a linear map from that lattice into \(A\) has
finite image and is continuous. Consequently the integral index \(\operatorname{len}(L_w/J_w)\) bounds the excess over the finite condition on the complementary row. The character rows require their own local conditions. \label{cor:split-selmer-4}
\end{proof}
\subsubsection*{Character row and column conditions}

The complementary row has now been accounted for. On the remaining rank-one blocks the cyclotomic twists distinguish the finite conditions, and the row and column impose different restrictions even at the same coefficient prime.

\begin{proposition}[Local conditions on the character blocks]
At a column the specialization is finite on block \(0\),
unramified on block \(1\), and zero on blocks \(k\ge2\).
At a row above \(\ell\), for \(1\le k\le t\) let \(L_k\) be the normalized length of the excess over the finite subgroup on block \(k\). Then
\[
 L_k\le([E_w:\mathbb Q_\ell]+2)c_B,\qquad c_B=40(t+1)(n+1)\lambda_N(h).
\]
\end{proposition}
\begin{proof}
At a column, the last defining relation of \(M\) reads in local
cohomology as a combination of the classes \(K_{u,a}\) and
\(K_{u,0}\) after evaluation. In block \(0\) the first classes
are finite and the last term vanishes because \(K_{u,0}\) has a
right factor \(E\). In block \(1\) the first terms vanish because
\(E_1\) acts by zero, and the last class is unramified. In
blocks \(k\ge2\) both sets of terms vanish. The generator
\(n_u\) contributes a coboundary in each case. These conclusions
hold at the normalized point and use the specific column lifts
of the determinant construction.

For the character rows, form the stable lattice \(Y\oplus O\)
with action
\[
  (z,a)\longmapsto (U_g z+a\,y(g),a).
\]
With \(e=(0,1)\), the quotient is the trivial line, and
\(ge-e=(y(g),0)\). Put \(i(z,a)=(z+aE,a)\).
The cocycle is \(y(g)=(U_g-1)E\), so
\[
    i(U_g z+a\,y(g),a)=(U_g(z+aE),a)=g\,i(z,a).
\]
The map \(i\) is injective, with the left \(U_g\)-action
on matrices and the trivial action on the second factor.

Evaluation
on the matrix columns embeds the rational space of this lattice
into a finite sum of the twisted family \(R_x^\circ\) and one
trivial line. The trivial character occurs in \(R_\pi^\circ\). The target of
the embedding is therefore crystalline, and so is its rational
subrepresentation \(V=(Y\oplus O)\otimes_O F\). For each
embedding \(\tau\), strictness of the crystalline filtration gives
\[
 \operatorname{Jumps}_\tau(V)
     \subseteq \bigcup_x \operatorname{Jumps}_\tau(R_x^\circ)
     \subseteq [-h,h].
\]
The union on the right over all embeddings has cardinality at
most \(n\) by the common row bound, with \(t\le h\).

The character lifting bound is used with that lattice, the
surjection onto the trivial line and the lift \(e\) just given.
It applies to an equivariant map of the kernel \(Y\) into the
divisible quotient of a rank-one lattice acted on by
\(\epsilon^{1-k}\), for \(1\le k\le t\). Extend this integral
map \(O\)-linearly to the full rational source using injectivity
of \(F/O\) over \(O\). The comparison only requires equivariance
on the kernel lattice.

Indeed, if \(\phi\) is the equivariant map on \(Y\) and
\(\widetilde\phi\) any such extension, then
\begin{equation}\label{eq:selmer-row-kernel-cocycle}
 c(g)=\widetilde\phi(ge-e)=\phi(y(g)),\qquad
 c(gh)=c(g)+\epsilon^{1-k}(g)c(h).
\end{equation}
Both identities in \eqref{eq:selmer-row-kernel-cocycle} use only the restriction to \(Y\). The completion lies in
the local level-\(\operatorname{lcm}(\ell,N)\) cyclotomic field,
the coefficients contain those roots of unity, and \(4\mid N\)
at \(\ell=2\). With
\[
 c_B=40(t+1)(n+1)\lambda_N(h),
\]
the
comparison places \(\ell^{c_B}[c]\) in the finite subgroup of
the target. For \(k=1\) this subgroup is unramified. For \(k>1\)
it is zero, by the rank-one jump \(k-1>0\) and the absence of
rational invariants in the Bloch--Kato dimension formula.
\cite{bloch-kato}\label{cor:split-selmer-5}

For the length allowance, let \(F_{k,w}\) be this finite
subgroup. It is divisible, as the image of a vector space.
Choose \(b\in F_{k,w}\) with \(\ell^{c_B}b=\ell^{c_B}[c]\).
Then \([c]-b\) represents the same class modulo \(F_{k,w}\) and
belongs to
\[
 H^1(E_w,A_k)[\ell^{c_B}]
   =\operatorname{im}\bigl(H^1(E_w,A_k[\ell^{c_B}])
                          \longrightarrow H^1(E_w,A_k)\bigr),
\]
by the cohomology sequence for multiplication by \(\ell^{c_B}\). For \(M_k=A_k[\ell^{c_B}]\), local Tate duality and the bound on
the coefficient module give \(\operatorname{len}H^i(E_w,M_k)\le c_B\)
for \(i=0,2\). The local Euler characteristic gives
\[
 \operatorname{len}H^1(E_w,M_k)
 = [E_w:\mathbb Q_\ell]c_B+
   \operatorname{len}H^0(E_w,M_k)+\operatorname{len}H^2(E_w,M_k).
\]
This gives an allowance for each character row even when a sharper row estimate is available. It leaves the inertia conditions away from the coefficient prime to be imposed.
\end{proof}
\subsubsection*{Tame tests on the character blocks}

It remains to control ramification away from the split row and column places. The tame inertia relations on each character block bound the quotient by unramified classes at the auxiliary level primes.

For a character block let \(A_k=(F/O)(1-k)\) and let
\(c(g)=\varphi(y(g))\) for an \(\mathcal A\)-linear map
from the column presentation into \(A_k\).

\begin{lemma}[Tame inertia and the remaining global condition]
Away from the split row and column places the character class is
unramified outside \(Q\); at \(v\in Q\) the ramified quotient \(D_v\) satisfies
\[
 \operatorname{len}(D_v)\le v_\ell(\zeta_v-1).
\] After imposing
the row, column and unramified conditions the residual length is at
most \(v_\ell(h_E)\) for \(k=1\); for \(k\ge2\) it is bounded by
the strict length with exponent \(k-1\) over the same field.
\end{lemma}
\begin{proof}
At a place
outside the split row and column places and the auxiliary
sets \(S,Q\), the family inertia acts trivially on every
fibre and \(\chi_1\) is trivial there. The last assertion
can be checked on its summand at \(\pi\). Thus \(y(g)=0\)
on this inertia.

At \(v\in Q\), the chosen inertial element
\(\gamma_v\) maps every fibre into the two eigenspaces for
\(1\) and \(\zeta_v=\chi_1(\gamma_v)\). The defining conditions on the tame generator give \(\zeta_v\ne1\), equality of all
\(\chi_k\) on inertia, and a common power \(n(g)\) of
\(R(\gamma_v)\) and of \(\chi_1(\gamma_v)\) for the actions
of each \(g\) on this inertia. The common power comes from
the finite image of the chosen tame character. In the family
algebra we have
\[
 (R(\gamma_v)-1)y(\gamma_v)
 =\zeta_v^{-1}(R(\gamma_v)-1)(R(\gamma_v)-\zeta_v)E=0.
\]

Evaluating under \(\varphi\) gives
\((\zeta_v-1)c(\gamma_v)=0\). On the target \(U_{\gamma_v}\) acts trivially. On ordinary powers the
cocycle therefore satisfies
\[
 c(\gamma_v^a)=\sum_{i=0}^{a-1}U_{\gamma_v}^ic(\gamma_v)=a\,c(\gamma_v)
 \qquad(a\ge1).
\]
Continuity on the procyclic quotient gives \(c(g)=n(g)c(\gamma_v)\).
Coboundaries vanish on inertia here. Therefore the quotient
by the unramified condition injects, by evaluation at
\(\gamma_v\), into \(A_k[\zeta_v-1]\), of normalized length
\(v_\ell(\zeta_v-1)\). This bound uses only \(\zeta_v-1\ne0\).

At \(v\in S\), there is an integral family operator \(a\)
such that
\[
 a(R(g)-\chi_1(g))E=0
\]
on this inertia and
\(a\) acts invertibly on \(A_k\). At the two type primes,
\(a=P_v(R(\gamma_v))\), where the coefficients of \(P_v\)
are in the localized Hecke algebra. The tame polynomial
condition asserts that on each fibre this polynomial kills
\(R(g)-\chi_1(g)\) and that \(P_v(\chi_1(\gamma_v))\) is
a unit upon evaluation at \(\pi\). Since all the scalar
characters agree on the chosen inertia, the second property
is exactly invertibility on \(A_k\).

At the remaining
\(S\)-places the family is simultaneously diagonal on inertia
with values \(1\) or \(\chi_1\) on the corresponding
summands. The factors kill the trivial and the \(\chi_1\) summands respectively. Hence on this inertia
\[
 (R(\gamma_v)-1)(R(g)-\chi_1(g))=0.
\]
At these places
\(\zeta_v-1\) is a unit, so the former operator is invertible
on the character target. Consequently \(c\) vanishes on
inertia at every \(S\)-place. These restriction bounds apply under their respective tame hypotheses.
\label{cor:split-selmer-6}

Subtract the row and \(Q\)-local allowance by taking the
kernel of restriction to the corresponding quotients. The
strict conditions for \(k\ge2\) at the coefficient prime
are zero on row and column decomposition groups. For \(k=1\)
the local finite condition for the trivial character at those
places is unramified. The declared rows and columns partition
the places over \(\ell\), and the remaining inertia groups
are among \(S,Q\) or the places with trivial family inertia.
In block \(k=1\), global class field theory therefore gives
\[
 \ker(\text{restriction to the local quotients})\hookrightarrow\operatorname{Hom}(\operatorname{Cl}(E),F/O).
\]
Its normalized length is at
most \(v_\ell(h_E)\).

For \(k\ge2\), let \(S_\ell\) consist of
the places over \(\ell\) and the infinite places. Put
\(A=(\mathbb Q_\ell/\mathbb Z_\ell)(1-k)\). The strict group is
\[
 S_{k-1}=\ker\!\left(H^1(G_{E,S_\ell},A)
        \longrightarrow\prod_{w\mid\ell}H^1(E_w,A)\right).
\]
The restriction kernel maps into \(S_{k-1}\otimes O\).
Indeed the local class vanishes at \(\ell\), all other places
are unramified, and finite freeness over \(\mathbb Z_\ell\) lets
tensoring by \(O\) commute with cochains and localization kernels.
Each \(\mathbb Z_\ell\)-length-one factor becomes \(O/\ell O\),
so the normalized length is \(\operatorname{length}_{\mathbb Z_\ell} S_{k-1}\). The strict
cohomology bound below applies with this exponent and on
this same auxiliary field. \label{cor:split-selmer-7}

\end{proof}
For block \(0\), the same vanishing at the places with trivial
family inertia follows directly from \(y(g)=0\). The row and
column conditions in the restriction kernel lie in the finite
condition on \(W/T\) at all places over \(\ell\). These are
the unramified and finite conditions of the relaxed Selmer
group at \(\pi\) used for \(s_\pi\). The length at this fixed point can therefore be compared with the signature Selmer length, provided the auxiliary character is the one used in its theta lattices.
